\documentclass[11pt,a4paper]{article}
\usepackage[T1]{fontenc}
\usepackage{lmodern,amsmath,amssymb,amsthm,mathtools,microtype}
\usepackage[margin=27mm]{geometry}
\usepackage[colorlinks=true,linkcolor=blue,citecolor=blue,urlcolor=blue]{hyperref}
\newtheorem{theorem}{Theorem}[section]
\newtheorem{lemma}[theorem]{Lemma}
\newtheorem{proposition}[theorem]{Proposition}
\newtheorem{corollary}[theorem]{Corollary}
\newcommand{\Q}{\mathbb Q}\newcommand{\Z}{\mathbb Z}
\newcommand{\OK}{\mathcal O_K}
\setlength{\emergencystretch}{2em}
\title{Power shifts and finite exceptions in Diophantine definitions}
\author{}\date{17 November 2026}
\begin{document}\maketitle
\begin{abstract}
A fixed integral power shift removes finite exceptions from an
integrality criterion without adding a quantified coordinate.
For a number field of degree $d$, a Diophantine set containing all
but finitely many rational integers and having only finitely many
nonintegral elements gives a definition of its ring of integers in
$dN$ total coordinates when the original definition uses $N$.
The resulting single equation has degree at most $2dD$ if the
original polynomial has degree at most $D$. A stronger finite-error
hypothesis gives a definition of the rational integers in the
original $N$ coordinates. Its degree is at most $2D$ whenever
the field has no quadratic subfield, in particular whenever its
degree is odd. Every shift may be constrained to a
prescribed integral arithmetic progression. No detector satisfying
these hypotheses over all number fields is constructed here.
\end{abstract}

\section{Avoiding finite sets by a power shift}
Throughout, $K$ is a number field of degree $d=[K:\Q]$, and all
Diophantine witnesses range over $K$. Coefficients may lie in $K$.
The total coordinate count includes the distinguished input.
For $m\ge2$ write $K^m=\{x^m:x\in K\}$.

\begin{lemma}\label{lem:avoid}
Let $T\subset K$ be finite, let $m\ge2$, and let $b_0\in\Z$ and
$M_0\ge1$. There are integers $b_1$ and $M_1\ge1$ such that
\[
 b_1+M_1\Z\subseteq b_0+M_0\Z,
 \qquad (K^m+B)\cap T=\varnothing
 \quad\text{for every }B\in b_1+M_1\Z.
\]
In particular the shifts can be chosen arbitrarily large and positive.
\end{lemma}
\begin{proof}
The assertion is immediate when $T$ is empty. Otherwise let $L/\Q$
be a finite Galois extension containing $K$ and all $m$th roots of
unity. Infinitely many rational primes split completely in $L$,
by the identity-class case of the Chebotarev density theorem
\cite[Theorem~8.31 and Corollary~8.32]{Milne}.
For each $\alpha\in T$ choose a distinct such prime $\ell_\alpha$
which divides neither $mM_0$ nor a denominator of $\alpha$.
Choose a prime $\mathfrak p_\alpha$ of $K$ above $\ell_\alpha$.
Its residue field is $\mathbb F_{\ell_\alpha}$, and
$\ell_\alpha\equiv1\pmod m$.

The cyclic group $\mathbb F_{\ell_\alpha}^{\times}$ has elements
which are not $m$th powers. Choose one, denoted $t_\alpha$, and
impose the rational integer congruence
\[
 B\equiv\overline\alpha-t_\alpha\pmod{\ell_\alpha},
\]
where the bar denotes reduction at $\mathfrak p_\alpha$.
Together with $B\equiv b_0\pmod{M_0}$ these congruences are
compatible by the Chinese remainder theorem. Let
$b_1+M_1\Z$ be their common solution progression.

If $x^m+B=\alpha$ for $x\in K$, then $x^m=\alpha-B$ is a
$\mathfrak p_\alpha$-adic unit. Hence $x$ is also a unit there,
and reduction gives $\overline x^{\,m}=t_\alpha$, a contradiction.
This proves the assertion for every shift in the progression.
\end{proof}

The construction is finite once $T$ is supplied as exact algebraic
numbers: split primes and the required residues can be found by
enumeration. If the finiteness of $T$ is known without an effective
description, the lemma supplies an existence result for the shift.
No effective determination of that set is asserted.

\section{Integrality criteria with the same witness count}
Suppose $S\subset K$ is defined by an integer $N\ge1$ and
\begin{equation}\label{eq:source}
 x\in S\quad\Longleftrightarrow\quad
 \exists\mathbf y\in K^{N-1}:P(x,\mathbf y)=0,
 \qquad P\in K[X,\mathbf Y],\quad\deg P\le D.
\end{equation}

\begin{theorem}\label{thm:ring}
Suppose both $\Z\setminus S$ and $S\setminus\OK$ are finite.
In every prescribed integral arithmetic progression there is a
subprogression of shifts $B$ for which
\begin{equation}\label{eq:intermediate}
 S_B=\{x\in K:\exists\mathbf y\in K^{N-1},
                      \ P(x^2+B,\mathbf y)=0\}
 \quad\text{satisfies}\quad \Z\subseteq S_B\subseteq\OK.
\end{equation}
Thus $S_B$ has a definition in $N$ total coordinates of degree at
most $2D$. The ring $\OK$ has a definition by $d$ equations in
$dN$ total coordinates, each of degree at most $2D$, and by one
equation in those same coordinates of degree at most $2dD$.
\end{theorem}
\begin{proof}
Apply Lemma~\ref{lem:avoid} with $m=2$ and
\[
 T=(S\setminus\OK)\cup(\Z\setminus S).
\]
For $n\in\Z$, the integer $n^2+B$ avoids $\Z\setminus S$, so
belongs to $S$. This proves $\Z\subseteq S_B$.
If $x\in S_B$, then $x^2+B\in S$ and avoids
$S\setminus\OK$, hence lies in $\OK$. Therefore
$x^2\in\OK$. The monic polynomial $Z^2-x^2\in\OK[Z]$ makes
$x$ integral over $\OK$, and the integral closedness of $\OK$
in $K$ gives $x\in\OK$. Equivalently, at every finite prime
$\mathfrak p$ the inequality $2v_\mathfrak p(x)\ge0$ implies
$v_\mathfrak p(x)\ge0$.

Choose an integral basis $\omega_1,\ldots,\omega_d$ of $\OK$.
Then
\[
                     \OK=\sum_{i=1}^d\omega_i S_B.
\]
For the forward inclusion expand an algebraic integer in this basis
with integer coefficients. For the reverse inclusion use
$S_B\subseteq\OK$.
Retain $s_1,\ldots,s_{d-1}$ and $d$ independent witness tuples
$\mathbf y_i\in K^{N-1}$, and put
\[
 s_d=\frac{X-\sum_{i<d}\omega_i s_i}{\omega_d}.
\]
The $d$ residual equations
\[
             R_i=P(s_i^2+B,\mathbf y_i)=0
             \qquad(1\le i\le d)
\]
use exactly
\[
                 1+(d-1)+d(N-1)=dN
\]
coordinates and have degree at most $2D$.
Choose a field extension $F/K$ of degree $d$ and a $K$-basis
$e_1,\ldots,e_d$; for $d=1$ take $F=K$. Such an extension exists
by an Eisenstein polynomial at a finite prime of $K$.
The single equation
\[
                N_{F/K}\left(\sum_{i=1}^d e_iR_i\right)=0
\]
is equivalent to all $R_i=0$, because the norm of a field element
vanishes exactly at zero. Its degree is at most $2dD$.
\end{proof}

\begin{corollary}\label{cor:positive-integrality}
Suppose $S\setminus\OK$ is finite and $S$ contains every
sufficiently large positive rational integer. In every prescribed
integral arithmetic progression, all sufficiently large positive
members of a suitable subprogression give
$\Z\subseteq S_B\subseteq\OK$ with the same square-shift
definition \eqref{eq:intermediate}. The bounds $dN$ and $2dD$
for a single equation defining $\OK$ remain valid.
\end{corollary}
\begin{proof}
Apply Lemma~\ref{lem:avoid} with $T=S\setminus\OK$ and $m=2$.
Choose $B$ above a threshold beyond which all positive integers
belong to $S$. Then $n^2+B\in S$ for every $n\in\Z$.
The integrality argument and the basis and norm construction in
Theorem~\ref{thm:ring} are unchanged.
\end{proof}

\begin{corollary}\label{cor:almost-ring}
If the symmetric difference of $S$ and $\OK$ is finite, then
$\OK$ itself has a definition in the same $N$ total coordinates
and degree at most $2D$.
\end{corollary}
\begin{proof}
In Lemma~\ref{lem:avoid} take $m=2$ and
$T=(S\setminus\OK)\cup(\OK\setminus S)$. For $x\in\OK$,
the integral element $x^2+B$ avoids the omitted set and belongs to
$S$. Conversely, the integrality argument in
Theorem~\ref{thm:ring} applies. Thus $S_B=\OK$.
\end{proof}

\section{The rational integers over a number field}
Say that an exponent $m\ge2$ \emph{reflects rationality in $K$}
if $x\in K$ and $x^m\in\Q$ imply $x\in\Q$.

\begin{lemma}\label{lem:power-rationality}
An exponent $m\ge2$ reflects rationality in $K$ if
\[
       \gcd(m,d)=1\quad\text{and}\quad \mu_m(K)\subseteq\Q,
\]
where $\mu_m(K)=\{\zeta\in K:\zeta^m=1\}$.
The exponent $2$ reflects rationality in $K$ if and only if
$K$ contains no quadratic subfield. In particular, it does so
whenever $d$ is odd.

Every number field has a rationality-reflecting exponent: an odd
prime $p$ with $p-1>d$ suffices. More generally, an odd prime
$p\nmid d$ suffices whenever $\mu_p(K)=\{1\}$. For a totally
real field the latter root-of-unity condition is automatic.

For any rationality-reflecting exponent $m$, the implication
$x^m\in\Z\Rightarrow x\in\Z$ holds for $x\in K$.
\end{lemma}
\begin{proof}
Suppose first that the two displayed conditions hold and that
$a=x^m\in\Q^\times$. Set $c=N_{K/\Q}(x)\in\Q^\times$.
Taking norms gives $c^m=a^d$. Choose integers $r,s$ with
$rm+sd=1$. Then
\[
                       a=(a^r c^s)^m.
\]
Thus $x/(a^rc^s)\in\mu_m(K)\subseteq\Q$, so $x\in\Q$.
The case $x=0$ is immediate.

If $x\notin\Q$ and $x^2\in\Q$, then $\Q(x)$ is a quadratic
subfield of $K$. Conversely, a quadratic subfield contains an
irrational square root of a rational number. This proves the
square criterion. An odd-degree extension cannot contain a
quadratic subfield, by the tower formula.

If $p-1>d$, a nontrivial $p$th root of unity would generate a
subfield of degree $p-1$, which is impossible. Also $p>d$, so
$p\nmid d$. For any odd prime, the only real $p$th root of
unity is one. These facts give the remaining sufficient criteria.
Finally a rational number whose positive integral power is an
integer is an integer, by reduction of numerator and denominator.
\end{proof}

\begin{theorem}\label{thm:Z}
Suppose the symmetric difference of $S$ and $\Z$ is finite.
Let $m\ge2$ reflect rationality in $K$. In every prescribed integral
arithmetic progression there is a subprogression of shifts $B$
such that
\[
 x\in\Z\quad\Longleftrightarrow\quad
 \exists\mathbf y\in K^{N-1}:P(x^m+B,\mathbf y)=0.
\]
This definition uses $N$ total coordinates and has degree at most
$mD$. If $K$ has no quadratic subfield, the degree bound is $2D$.
\end{theorem}
\begin{proof}
Apply Lemma~\ref{lem:avoid} with exponent $m$ and
$T=(S\setminus\Z)\cup(\Z\setminus S)$. If $x\in\Z$, then
$x^m+B$ is an integer avoiding $\Z\setminus S$, so it belongs
to $S$. Conversely, membership in $S$ and avoidance of
$S\setminus\Z$ imply $x^m+B\in\Z$, hence $x\in\Z$ by
Lemma~\ref{lem:power-rationality}.
\end{proof}

\begin{corollary}\label{cor:positive-Z}
Suppose $S\setminus\Z$ is finite and $S$ contains every
sufficiently large positive rational integer. Let $m\ge2$
reflect rationality in $K$. In every prescribed integral arithmetic
progression there are arbitrarily large shifts $B$ such that
\[
 x\in\Z\quad\Longleftrightarrow\quad
 \exists\mathbf y\in K^{N-1}:\quad
 \begin{cases}
 P(x^m+B,\mathbf y)=0,&m\text{ even},\\
 P(x^m+B,\mathbf y)P(-x^m+B,\mathbf y)=0,&m\text{ odd}.
 \end{cases}
\]
This definition uses the same $N$ total coordinates and has
degree at most $mD$ for even $m$ and $2mD$ for odd $m$.
In particular, when $K$ has no quadratic subfield a single
square-shift equation suffices, with degree at most $2D$.
\end{corollary}
\begin{proof}
Apply Lemma~\ref{lem:avoid} with $T=S\setminus\Z$ and
exponent $m$, and choose $B$ above the positive completeness
threshold. When $m$ is even, $n^m+B\ge B$ for all $n\in\Z$,
so the single equation is complete. Its soundness follows as in
Theorem~\ref{thm:Z}.

When $m$ is odd, both $x^m+B$ and $-x^m+B=(-x)^m+B$ avoid $T$.
If the displayed product vanishes, one factor vanishes, so one
of these two inputs is an integer. Lemma~\ref{lem:power-rationality}
then gives $x\in\Z$. Conversely, if $x\in\Z$, at least one
of the two inputs is an integer at least $B$, hence belongs to
$S$. A witness tuple for that input makes the product vanish.
The same tuple can be used in both factors because only one
factor is required to vanish.
\end{proof}

For example, a real quadratic field admits the exponent $m=3$,
giving degree bounds $3D$ in Theorem~\ref{thm:Z} and $6D$ in
Corollary~\ref{cor:positive-Z}. More generally, a totally real
field admits every odd prime not dividing its degree. In
$\Q(\sqrt{-3})$, the cube does not reflect rationality because
of the nontrivial cube roots of unity, whereas the exponent $5$
satisfies Lemma~\ref{lem:power-rationality}.

\section{Height hypotheses and limits of the conclusions}
\begin{corollary}\label{cor:height}
The conclusions of Theorem~\ref{thm:ring} hold if
$\Z\setminus S$ is finite and the elements of $S\setminus\OK$
have bounded absolute logarithmic height. The conclusions of
Theorem~\ref{thm:Z} hold if $\Z\setminus S$ is finite and the
elements of $S\setminus\Z$ have bounded height.
\end{corollary}
\begin{proof}
Northcott's theorem makes a bounded-height subset of the fixed
number field $K$ finite \cite[Part~B, Section~2]{HS}.
For completeness, if $\alpha$ has degree $e\le d$ and height
at most $C$, its primitive integral minimal polynomial has
Mahler measure $\exp(eh(\alpha))\le\exp(dC)$.
Every coefficient then has absolute value at most
$2^d\exp(dC)$, by the elementary symmetric function formulas.
Only finitely many such integral polynomials exist.
Apply the corresponding theorem.
\end{proof}

All substitutions preserve the number of existential witnesses.
They change the accepted parameter set by composition; they do not
assert equality with the set obtained by deleting and adjoining
arbitrary finite subsets of $S$. The coefficients of the shifts
may be ineffective when only finiteness of the exceptional sets is
known. The bounds concern degree and coordinates, and supply no
uniform coefficient-height estimate in that situation.

In particular, a height inequality for elliptic points alone does
not furnish the Diophantine set required in
Corollary~\ref{cor:height}. Its integer completeness and the height
bound on false inputs remain separate hypotheses. No assertion
that $\OK$ is Diophantine in every number field follows without
constructing such a set.

Uniform bounds for false inputs arise from the power-law and
logarithmic height comparisons used in rational integer
certificates; see \cite[Sections~7 and~11]{Sun}.
When such a bound is available for a complete source, the
substitution above removes its finite false-input set.
This use of a fixed power shift is distinct from the computation
of a fixed multiplication map in \cite[Section~9]{Sun}.

\begin{thebibliography}{9}
\bibitem{Milne} J. S. Milne, \emph{Algebraic Number Theory},
version 3.08, 19 July 2020, Theorem 8.31 and Corollary 8.32.
\url{https://www.jmilne.org/math/CourseNotes/ANT.pdf}.
\bibitem{HS} M. Hindry and J. H. Silverman,
\emph{Diophantine Geometry: An Introduction},
Graduate Texts in Mathematics 201, Springer, 2000,
Part B, Section 2. \url{https://doi.org/10.1007/978-1-4612-1210-2}.
\bibitem{Sun} Z.-W. Sun, \emph{$\mathbb Z$ is diophantine over
$\mathbb Q$ with 26 unknowns}, arXiv:2611.14605v2,
13 November 2026. \url{https://arxiv.org/abs/2611.14605v2}.
\end{thebibliography}
\end{document}
